\documentclass[conference]{IEEEtran}
\usepackage{amsmath,amssymb}
\title{Migrating KV Cache under Bandwidth Uncertainty in Edge Collaboration}
\author{Author A}
\begin{document}
\maketitle

\begin{abstract}
Collaborative edge inference can migrate key-value cache between nodes to balance memory pressure.
Migration decisions are difficult because link bandwidth fluctuates and cache state is large.
We model migration as a sequential decision problem with partial observations and propose a threshold policy that triggers migration only when the expected benefit exceeds a bandwidth-dependent cost.
On a five-node testbed, the policy reduces preemption by 31\% while keeping migration traffic below 9\% of the total uplink budget.
\end{abstract}

\section{Introduction}
\label{sec:intro}
Key-value cache has become the dominant memory cost of long-context inference.
When a session moves between edge nodes, the cache must either be rebuilt or transferred.
Rebuilding wastes computation; transferring competes with request traffic.
This paper studies when migration is worth its bandwidth cost.
We argue that a fixed threshold on cache size is insufficient, because the same transfer can be cheap on one link and prohibitive on another.
Our contributions are a partial-observation model of bandwidth, a threshold policy with a provable regret bound, and an empirical study of migration under contention.

\section{Related Work}
\label{sec:related}
Cache migration appears in serverless inference and in distributed serving systems.
Prior systems migrate eagerly to minimize recomputation, or lazily to avoid bandwidth spikes.
Neither approach conditions the decision on the observed distribution of available bandwidth.
Our formulation treats bandwidth as a latent state and updates a belief after each transfer.

\section{Problem Formulation}
\label{sec:problem}
Let $c_t$ denote the cache size at epoch $t$ and let $b_t$ denote the available bandwidth.
Migration costs $\alpha c_t / b_t$ seconds and saves expected recomputation $r(c_t)$.
The controller observes a noisy signal $y_t = b_t + \varepsilon_t$ before deciding.
We seek a policy that minimizes the sum of migration and recomputation costs over a horizon $T$.

\section{Method}
\label{sec:method}
We maintain a belief over $b_t$ and trigger migration when
\begin{equation}
\label{eq:threshold}
r(c_t) - \alpha c_t / \hat{b}_t > \tau,
\end{equation}
where $\hat{b}_t$ is the posterior mean and $\tau$ is a hysteresis margin.
The margin prevents oscillation when the belief is uncertain.
We show that the resulting policy has sublinear regret relative to an oracle that knows $b_t$.

\section{Evaluation}
\label{sec:eval}
We emulate five edge nodes with shaped links and long-context sessions.
The threshold policy reduces preemption by 31\% compared with eager migration.
Migration traffic stays below 9\% of the uplink budget.
In particular, the hysteresis margin accounts for most of the gain under bursty bandwidth.
A sensitivity study shows that the policy degrades gracefully when the signal noise doubles.

\section{Conclusion}
\label{sec:conclusion}
We studied cache migration under bandwidth uncertainty.
A belief-based threshold policy avoids both eager transfers and costly recomputation.
Future work will consider multi-hop migration and cache compression.

\end{document}
